% Generated by scripts/gen_blueprint.py from blueprint/nodes/11-dilogarithm-dichotomy.toml. Do not edit.

\chapter{A dichotomy for the dilogarithm at one half}
\label{chap:11-dilogarithm-dichotomy}

Two statements about classical constants from the library's work on Diaz's conjecture. Both rest on the four exponentials theorem in transcendence degree one (Chapter~\ref{chap:06-four-exponentials}). Each alternative of Corollary~\ref{cor:dilog_half_irrational_or_exp_i_div_pi_transcendental} is open, so at least one of two open questions has a positive answer.

% Prove2Me: DiazModulus.recip_pi_log_of_rational_quadratic_relation -- https://prove2.me/theorems/de445e9d-80c9-4aef-864b-71b412e9ea1a
\begin{proposition}[A rational quadratic relation between $t$ and $\pi$]
\label{prop:recip_pi_log_of_rational_quadratic_relation}
\lean{Diaz.recip_pi_log_of_rational_quadratic_relation}
\leanok
Let $t\neq0$ be real with $e^{t}$ algebraic, and let $a,b,c\in\Q$ with $a\neq0$ and $at^{2}+b\pi^{2}=c$. Then $e^{i\gamma/\pi}$ is transcendental for every $\gamma\in\Q$, $\gamma\neq0$.

{\small\emph{Source:} Brownawell's Corollary~5 \cite[p.~23]{Bro74} at $\eta=t/\pi$, followed by rational scaling. The statement for general $a,b$ was not found in the sources read.}
\end{proposition}

\begin{proof}
\leanok
\uses{thm:four_exponentials_trdeg_one,lem:trdeg_adjoin_le_one_of_isAlgebraic_adjoin}
Suppose $e^{i\gamma/\pi}$ is algebraic, so that $\lambda=i\gamma/\pi$ is a logarithm of an algebraic number. The relation gives $t^{2}/(i\pi)=-\frac{c}{a\gamma}\,\lambda+\frac{b}{a}\,i\pi$, again a logarithm of an algebraic number. The matrix $\begin{pmatrix}t&t^{2}/(i\pi)\\ i\pi&t\end{pmatrix}$ has determinant $0$ and entries in a field of transcendence degree one, since $t$ is algebraic over $\Q(\pi)$. By Theorem~\ref{thm:four_exponentials_trdeg_one} its rows or its columns are linearly dependent over $\Q$; either way $\alpha t+\beta i\pi=0$ with $\alpha,\beta\in\Q$ not both zero, impossible since $t$ is real and non-zero and $i\pi$ is purely imaginary.
\end{proof}

% Prove2Me: DiazModulus.dilog_half_irrational_or_exp_i_div_pi_transcendental -- https://prove2.me/theorems/f9999428-837c-4d87-a70f-5790d5c414eb
\begin{corollary}[$\mathrm{Li}_2(1/2)$ against $e^{i/\pi}$]
\label{cor:dilog_half_irrational_or_exp_i_div_pi_transcendental}
\lean{Diaz.dilog_half_irrational_or_exp_i_div_pi_transcendental}
\leanok
Either $\pi^{2}/12-(\log2)^{2}/2$ is irrational, or $e^{i\gamma/\pi}$ is transcendental for every $\gamma\in\Q$, $\gamma\neq0$.

{\small\emph{Source:} By Euler's identity, $\pi^{2}/12-(\log2)^{2}/2=\mathrm{Li}_2(1/2)=\sum_{n\ge1}1/(n^{2}2^{n})$. Its irrationality is listed as unknown in \cite[p.~274]{Wal04} and is still open: \cite{Hat93,RV05,CDT24} cover $\mathrm{Li}_2(1/q)$ only for $q\ge6$ and $q\le-5$. The dichotomy was not found in the sources read.}
\end{corollary}

\begin{proof}
\leanok
\uses{prop:recip_pi_log_of_rational_quadratic_relation}
If $\pi^{2}/12-(\log2)^{2}/2=s\in\Q$, apply Proposition~\ref{prop:recip_pi_log_of_rational_quadratic_relation} with $t=\log2$, $a=-1/2$, $b=1/12$ and $c=s$.
\end{proof}
